\documentclass[conference]{IEEEtran}
\IEEEoverridecommandlockouts
\usepackage{cite}
\usepackage{amsmath,amssymb}
\usepackage{graphicx}
\usepackage{textcomp}
\usepackage{xcolor}
\usepackage{url}
\def\UrlFont{\footnotesize\ttfamily}
\def\BibTeX{{\rm B\kern-.05em{\sc i\kern-.025em b}\kern-.08em T\kern-.1667em\lower.7ex\hbox{E}\kern-.125emX}}

\begin{document}

\title{Federated Intrusion Detection Under Label-Flipping Poisoning:\\
A Simulated Multi-Site Study on CICIoT2023}

\author{
\IEEEauthorblockN{Salvador Leon Lopez, Elena Victoria Sanchez, Joe Reyna, April Flores}
\IEEEauthorblockA{Department of Computer Science\\
The University of Texas Rio Grande Valley\\
CSCI 6330: Foundations of Intelligent Security Systems, Fall 2026}
}

\maketitle

\section{Problem and Motivation}
Retail headquarters, hospitals and bank branches each see attack traffic the
others never do, yet privacy policy bars pooling raw captures.
Federated learning (FL) lets each site train locally and share only model
updates \cite{mcmahan2017}. Since the server never inspects the underlying
records, one compromised participant can corrupt the shared model
\cite{tolpegin2020}. We quantify that exposure for intrusion detection.

\section{Research Question and Expectation}
\textbf{RQ:} In a six-client federated IDS trained on non-IID partitions of
CICIoT2023, how much global detection performance is lost when one client
performs a label-flipping attack, and does robust aggregation recover it?

\textbf{Threat model and expectation:} One client relabels 50\% of its attack
samples as benign each round; all other parties are honest. Under FedAvg we expect
global macro F1 to fall by more than 3 points versus the clean run, and
trimmed-mean aggregation \cite{yin2018} to recover at least half of that loss;
smaller effects would falsify this.

\section{Dataset and Client Partition}
We use CICIoT2023 \cite{neto2023}, labeled IoT flows from the Canadian Institute
for Cybersecurity. Its official page states no license and asks for citation, so
we use it for coursework research only (accessed Sep.~20, 2026). It is
lab-generated single-testbed data with heavy class imbalance, so the six branch
offices are \textbf{simulated clients} produced by partitioning, and we make no
real-world claims. We build a documented stratified subset of about 300{,}000
flows, collapsing the 33 attack types into the seven published families plus
benign. A stratified 20\% test set and 10\% validation set are held out first (seed 42); the
test set is for final reporting only, and all settings are chosen on validation.
The other 70\% is grouped into six clients by k-means ($k=6$, seed 42) on
standardized flow features, fit on this pool only, so no labels build the
partition; an IID split is the control. Clients are never defined by labels, IPs,
ports, timestamps or row indices, and identifier fields (if present) are dropped
to prevent shortcut learning. We quantify leakage risk with per-client label
distributions and client--label mutual information.

\section{Method and Experimental Conditions}
Python, PyTorch and Flower \cite{beutel2020} run in single-machine simulation. One
small MLP over numeric flow features is held constant; training runs 30 rounds
of one local epoch, all six clients participating. Core configurations: (1) a
\textbf{centralized reference} on pooled data; (2) a \textbf{clean federated
baseline}, FedAvg on non-IID and IID splits; (3) a \textbf{poisoned condition} with
one label-flipping client; and (4) \textbf{robust aggregation}, trimmed mean, on the
poisoned and clean non-IID setups, since honest but heterogeneous clients can look
like outliers. Local-only models and median aggregation are extensions.

\section{Evaluation Plan}
The primary metric is macro F1 on the global test set, since accuracy alone is
insufficient on imbalanced data, with per-class precision, recall and FNR, FPR,
PR-AUC and confusion matrices. Runs use 3 seeds (mean and standard deviation). Federated runs add per-round performance, client-level dispersion, a
communication proxy (bytes per round) and runtime.

\section{Safety, Reproducibility and Timeline}
All experiments are simulated on the project's own static data and code; no real
service, traffic generation, university or third-party system, or personal data is
involved, and poisoning is a controlled label change compared against a clean
baseline. The repository will include pinned dependencies, fixed seeds, subset and
partition scripts, and a README reproducing every reported number. Milestone~1
(Oct.~25): subset, partitions, centralized model; Milestone~2 (Nov.~15): clean vs.\
poisoned runs; final (Nov.~30, Dec.~7): limitations, demo, defenses. Members own
data preparation, the centralized model, the federated baseline, and attack and
defense.

\footnotesize
\begin{thebibliography}{00}
\itemsep 0pt
\bibitem{mcmahan2017} B. McMahan \emph{et al.}, ``Communication-efficient learning of deep networks from decentralized data,'' \emph{Proc. AISTATS}, 2017.
\bibitem{tolpegin2020} V. Tolpegin, S. Truex, M. E. Gursoy, and L. Liu, ``Data poisoning attacks against federated learning systems,'' \emph{Proc. ESORICS}, 2020.
\bibitem{yin2018} D. Yin, Y. Chen, K. Ramchandran, and P. Bartlett, ``Byzantine-robust distributed learning,'' \emph{Proc. ICML}, 2018.
\bibitem{neto2023} E. C. P. Neto \emph{et al.}, ``CICIoT2023: A real-time dataset and benchmark for large-scale attacks in IoT environment,'' \emph{Sensors}, vol.~23, no.~13, 2023. \url{https://www.unb.ca/cic/datasets/iotdataset-2023.html}
\bibitem{beutel2020} D. J. Beutel \emph{et al.}, ``Flower: A friendly federated learning framework,'' \emph{arXiv:2007.14390}, 2020.
\end{thebibliography}

\end{document}
